


% ------------------ DOCUMENT SETUP ------------------ 
% The document class defines the document type (report) and sets the font size (10pt)
\documentclass[10pt]{report}
\author{Felipe Santos Almeida}

% Inputs the Document Packages
% ------------------ PACKAGES ------------------ 
% Packages add extra commands and features to your LaTeX document. 
% In here, some of the most common packages for a thesis document have been added 

% LaTeX's float package
\usepackage{float}

% LaTeX's color package
\usepackage{color}

% LaTeX's main math package
\usepackage{amsmath}

% LaTeX's Caption and subcaption packages
\usepackage[format=hang,font=normalsize,labelfont=bf,labelsep=colon,singlelinecheck=off]{caption}
\usepackage{subcaption}

% The graphicx package provides graphics support for adding pictures.
\usepackage{graphicx}

% Longtable allows you to write tables that continue to the next page.
\usepackage{longtable}

% The geometry packages defines the page layout (page dimensions, margins, etc)
\usepackage[a4 paper, top=25mm, bottom=25mm]{geometry}

% Defines the Font of the document, e.g. Arimo font (Check Fonts here: https://tug.org/FontCatalogue/)
\usepackage[sfdefault]{arimo}

% Font encoding
\usepackage[T1]{fontenc}

% This package allows the user to specify the input encoding
\usepackage[utf8]{inputenc}

% This package allows you to add empty pages
\usepackage{emptypage}

% Allows inputs to be imported from a directory
\usepackage{import}

% Provides control over the typography of the Table of Contents, List of Figures and List of Tables
\usepackage{tocloft}

% The setspace package controls the line spacing properties.
\usepackage{setspace}

% Allows the customization of Latex's title styles
\usepackage{titlesec}

% Allows the customization of Latex's table of contents title styles
\usepackage{titletoc}

% The package provides functions that offer alternative ways of implementing some LATEX kernel commands
\usepackage{etoolbox}

% Provides extensive facilities for constructing and controlling headers and footers
\usepackage{fancyhdr} 

% Typographical extensions, namely character protrusion, font expansion, adjustment 
%of interword spacing and additional kerning
\usepackage{microtype}

% Manages hyperlinks 
\usepackage[colorlinks=true,linkcolor=black,urlcolor=blue]{hyperref}

% Generates PDF bookmarks
\usepackage{bookmark}

% Add color to Tables
\usepackage[table,xcdraw]{xcolor}

% Use these two packages together -- they define symbols
%  for e.g. units that you can use in both text and math mode.
\usepackage{gensymb}
\usepackage{textcomp}
% Bibliography package
%\usepackage[backend=biber,style=nature]{biblatex}
%\addbibresource{references.bib} % Add the .bib file that contains the references


\usepackage[utf8]{inputenc}
\usepackage{amsmath}
\usepackage{csquotes}% Recommended

\usepackage[style=authoryear,open={(},close={)}]{biblatex}
\addbibresource{References.bib}

%\usepackage{natbib}
%\bibliographystyle{abbrvnat}
%\hypersetup{citecolor=black}
%\setcitestyle{authoryear,open={(},close={)}}


%\addbibresource{References.bib}% Syntax for version >= 1.2



% The natbib package allows book-type citations (e.g. (Saucer et al., 1993))
% More details are here: http://merkel.zoneo.net/Latex/natbib.php
%\usepackage[round]{natbib}

% The linenumbers command adds line numbers next to your text 
% That can be useful for discussing edits in drafts.
%\usepackage[modulo]{lineno}
%\linenumbers 


% This package provides an easy way to input latin sample text (for the template only)
\usepackage{blindtext}


\usepackage{booktabs}

% Controls how many subsections the document can take
%  and how many of those will get put into the contents pages.
\setcounter{secnumdepth}{2}
\setcounter{tocdepth}{2}

% Line Spacing
\setstretch{1.0}

% The folder path where the images will be uploaded
\graphicspath{{./Images/}} 

% Places a dot after Chapter/Section/Subsection number in Table of Contents
\renewcommand{\cftchapaftersnum}{.}
\renewcommand{\cftsecaftersnum}{.}
\renewcommand{\cftsubsecaftersnum}{.}

%  Customize Dot spacing in Table of Contents/List of Figures/Tables
\renewcommand{\cftdotsep}{0.3}

% Numeration Type for Chapters and Sections (Roman I, II, II / Arabic 1, 2, 3)
\renewcommand\thechapter{\Roman{chapter}}
\renewcommand\thesection{\arabic{section}}

% Line Break Properties
\tolerance=1
\emergencystretch=\maxdimen
\hyphenpenalty=10000
\hbadness=10000


% Formatting Table of Contents/Lists titles
\renewcommand{\contentsname}{\normalfont\bfseries\LARGE{CONTENTS}}
\renewcommand{\listfigurename}{\normalfont\bfseries\LARGE{LIST OF FIGURES}}
\renewcommand{\listtablename}{\normalfont\bfseries\LARGE{LIST OF TABLES}}


% Title Formatting customization
\titleformat{\chapter}{\normalfont\bfseries\LARGE}{\thechapter.}{1em}{\MakeUppercase}

\titleformat{\section}{\normalfont\bfseries\large}{\thesection.}{1em}{\MakeUppercase}
\titlespacing*{\section} {0pt} {0pt} {15pt} % left, before, after

\titleformat{\subsection}{\normalfont\bfseries\large}{\thesubsection.}{1em}{}
\titlespacing*{\subsection} {0pt} {10pt} {10pt}

\titleformat{\subsubsection}{\normalfont\bfseries\large}{}{1em}{}
\titlespacing*{\subsubsection} {0pt} {10pt} {10pt}


% HEADER AND FOOTER
\pagestyle{fancy}  % Set Page Style (Header and Footer Style)
\fancyhf{}  % Clears the header and footer (from the default info)

% Header
\renewcommand{\headrulewidth}{0pt}  % Removes the default Horizontal Line in Header
\fancyhead[L]{SN21125032}
\fancyhead[R]{May 2023}

% Footer
\fancyfoot[C]{\thepage} % Page Number

% Change figure numbering per section
\numberwithin{figure}{section}
\numberwithin{table}{section}



 \captionsetup{justification=centering}


%  -------------------------------------------------
%  --------- The document starts from here --------- 
%  -------------------------------------------------

\begin{document}

% -------------------  INTRODUCTION  ---------------------

\begin{center}
    % UCL IMAGE
    \vspace*{-3cm}
    \makebox[\textwidth]{\includegraphics[width=\paperwidth]{Images/UCL_LOGO_new.png}}
\end{center}   
    % Title
    {\LARGE\textbf{Coursework 2\\
    % Subtitle
    CASA0011 – Agent-Based Modelling\\}}
SN21125032 | Word count:

\vspace{5mm} %5mm vertical space
  
\section{Research Question}

\section{ODD Description}
\subsection{Purpose and patterns}

According to ..... This a serious problem related to the lack of street space for pedestrians, and how this affects the rate of injuries caused by car crashes. Models developed on Netlogo related to this field, such as 15min cities ( brief explanation), Town - Traffic and Crowd Simulation( brief explanation), Traffic grid( brief explanation) and Taxi Cabs( brief explanation) illustrate the potentials of micro-simulation related to this field.
Therefore, this model aims to test if a micro-simulation of a street intersection using pedestrian flows in crossing a street can explain the relationship between sidewalk dimension and waiting time to cross the street. In this model, it is applied...... However, this model does not intend to simulate a realistic street intersection with all the agents. The aim is to intend pedestrian behaviour in the context of commute flows and their relationship with the available space for a given sidewalk.  

Thus, the aim is to .....

The patterns, 

First, a high concentration of people close by pedestrian crossing
Second, alternation in the density of people in opposite regions of the intersection
Third, travessy time is not proportional to the concentration of people

\subsection{Entities, state variables, and scales}

\textbf{Entities}

The entities in the project are the pedestrians who are walking in random directions, aiming to cross the street to continue the walk. In that context, the model also has the street, the sidewalk and the building representing the built environment. Besides that, it also has a pedestrian crossing and a separation in the middle of the streets.


\textbf{State variables}


\begin{table}[h]
\centering
\begin{tabular}{@{}lll@{}}
\toprule
\textbf{Variable name} & \textbf{Variable type and units} & \textbf{Meaning}                                                                                         \\ \midrule
meaning                & String                           & Represents the role of the patch                                                                         \\ \midrule
will-cross?            & Boolean                          & \begin{tabular}[c]{@{}l@{}}Indicates whether anybody is going \\ to cross the street\end{tabular}        \\ \midrule
used                   & Numeric                          & \begin{tabular}[c]{@{}l@{}}How many  pedestrians are \\ using the crossing\end{tabular}                  \\ \midrule
speed                  & Numeric                          & The current speed of a person                                                                            \\ \midrule
walk-time              & Numeric                          & \begin{tabular}[c]{@{}l@{}}How long the person will walk \\ before crossing the road\end{tabular}        \\ \midrule
crossing-part          & Numeric                          & Divides the crossing into parts                                                                          \\ \midrule
waiting?               & Boolean                          & \begin{tabular}[c]{@{}l@{}}Indicates whether the pedestrian is \\ waiting to cross the road\end{tabular} \\ \midrule
num-of-people          & Numeric                          & \begin{tabular}[c]{@{}l@{}}The desired number of people \\ in the simulation\end{tabular}                \\ \midrule
time-to-crossing       & Numeric                          & \begin{tabular}[c]{@{}l@{}}Represents the time taken \\ to cross the road\end{tabular}                   \\ \bottomrule
\end{tabular}
\caption{State variables}
\label{table: variables}
\end{table}

\newpage

\textbf{Scales}

The model has not a specified scale. The size of the patches is not defined.

\subsection{Process overview and scheduling}


The model is developed to simulate pedestrians' movement when they cross the street. The pedestrians follow specific rules such as walking on the sidewalk, waiting at the wait point close to the pedestrian crossing and then crossing the streets. Thus, the micro-simulation progress each interaction in the main loop following this step, allowing for tracking the movement and control of the interaction through the number of pedestrians, time to crossing and speed. To summarize, the simulations occur in that sequence:

1. Setup the environment
2. Main loop
3. Move pedestrians
4. Walking rules
5. Crossing the street
6. Simulation track


\subsection{Design concepts}


\subsection{Input data}


\subsection{Submodels}





\section{Brief Methodology }



Limitations

Adding the flow of cars, cyclists and traffic lights would be relevant.

















% -------------------  BIBLIOGRAPHY ---------------------
\newpage
\printbibliography[title = {References}]
\addcontentsline{toc}{chapter}{References} % Adds References Section to Table of Contents

\vspace{5mm} %5mm vertical space

To access GitHub Repository \href{https://github.com/brfelipealmeida/ABM_coursework}{Link}

\end{document}
%  -------------------------------------------------
%  --------- The document ends from here ----------- 
%  -------------------------------------------------